% =================================================================================
% Hier ausw�hlen, ob TUD-Design oder nicht
% =================================================================================
\newif\ifTUDdesign
\TUDdesigntrue					% TUD-Design
%\TUDdesignfalse				% F�r Rechner ohne installierte TUDdesign-Pakete
% =================================================================================


% =================================================================================
% Hier Daten f�r studentische Arbeit eingeben
% =================================================================================
\newcommand{\SADATyp}{Praktikumsbericht}
\newcommand{\SADATitel}{Versuch II: Linearisierung, Steuerbarkeit und Beobachtbarkeit}
\newcommand{\SADAStadt}{Darmstadt}
\newcommand{\SADAAutor}{Andreas Jentsch, Ali Kerem Sacakli}
\newcommand{\SADABetreuer}{}
\newcommand{\SADABetreuerII}{}
\newcommand{\SADABetreuerIII}{}
\newcommand{\SADABegin}{}
\newcommand{\SADAAbgabe}{Praktikum Matlab/Simulink II}
\newcommand{\SADASeminar}{}
% =================================================================================


% =================================================================================
% Auswahl des IAT-Fachgebiets (rtm / rtp)
% =================================================================================
\newif\ifrtm
\rtmtrue	% rtm
%\rtmfalse	% rtp
% =================================================================================


% =================================================================================
% Erkl�rung, dass die Arbeit ohne Hilfe Dritter etc. erstellt wurde
% =================================================================================
\def\SADAVarianteErklaerung{ETIT}		% FB 18, Elektrotechnik
%\def\SADAVarianteErklaerung{MBDA}		% FB 16, Maschinenbau, Diplomarbeit
%\def\SADAVarianteErklaerung{MBSA}		% FB 16, Maschinenbau, Studienarbeit
% =================================================================================


% =================================================================================
% Ausnahmen von der automatischen Silbentrennung
% =================================================================================
\hyphenation{Aktu-ali-sie-rung Screen-shots Pa-rallel-ro-bo-ter Zu-stands-raum-mo-del-le nach-voll-zieh-bar Pro-jekt-se-mi-nar}
% =================================================================================


% =================================================================================
% Hier NIICHTS �ndern!
% =================================================================================
\ifTUDdesign
	\documentclass[11pt, twoside, colorback, accentcolor=tud2c, nopartpage, bigchapter, fleqn, ngerman, longdoc]{tudreport}
\else
	\documentclass[11pt, a4paper, twoside, fleqn, ngerman]{scrreprt}
  % F�r Entwurf auf Rechnern ohne installierte TUDdesign-Pakete	
	\usepackage{exscale}	% Korrektur math-Zeichen
	\usepackage{eurosym}
\fi
%eps-plot von Matlab in pdf umwandeln, um diese einbinden zu k�nnen
\usepackage{epstopdf}
%Paket listings f�r das Einbinden von Quelltext
\usepackage{listings}
%Stil Matlab_colored definieren
\lstdefinestyle{Matlab_colored}
{
	language = Matlab,
	tabsize = 4,
	framesep = 3mm,
	frame = tb,
	classoffset = 0,
	basicstyle = \ttfamily,
	keywordstyle = \bfseries\color[rgb]{0,0,1},
	commentstyle = \itshape\color[rgb]{0.133,0.545,0.133},
	stringstyle = \color[rgb]{0.672,0.126,0.941},
	extendedchars = true,
	breaklines = true,
	prebreak = \textrightarrow,
	postbreak = \textleftarrow,
	numbers = left,
	numberstyle = \tiny,
	stepnumber = 5
}


% Dieses File dient zum einbinden wichtiger und n�tzlicher Pakete.
% Nicht alle Pakete m�ssen verwendet werden.
%
\usepackage{t1enc}			% evtl. dc-Fonts
\usepackage[T1]{fontenc}	% F�r Silbentrennung bei W�rten mit Sonderzeichen (z.B. Umlaute)
\usepackage[latin1]{inputenc}
							% Um Sonderzeichen (�, �, �, ...) direkt eingeben zu k�nnen
\usepackage[english,ngerman]{babel}
							% F�r Sprachenspezifisches
							% ngerman ist schon als globale Option definiert

%\usepackage{helvet}			% Helvetica als Standard-Sans-Schriftart
\usepackage[stable]{footmisc}
\usepackage{booktabs}


\usepackage{graphicx}		% zum Einbinden von Postscript
\usepackage{psfrag}			% Beschriftung der Bilder
\usepackage{amsmath}		% Mehr mathematischen Formelsatz
%\usepackage{amssymb}		% Mehr mathematische Symbole
%\usepackage{amsthm}

\usepackage{float}			% F�r Parameter [H] bei Flie�objekten

\usepackage{epsfig}			% Um eps-Bilder einzubinden
\usepackage{scrhack}    % Um Warnung "float@addtolists detected" zu unterdr�cken 
\usepackage{subfig}			% F�r Unterabbildungen
\usepackage{ltxtable} 		% Vereinigt TabularX und Longtable
\usepackage{rotating}		% Zum Drehen von Objekten
\usepackage{bibgerm}		% F�r deutsche Literaturverwaltung
%\usepackage{wrapfig}		% F�r kleine Bilder am Rand
%\usepackage{floatflt}		% Alternative zu wrapfig
%\usepackage[hang]{caption}	% Damit mehrzeilige Bildunterschriften gut aussehen
\usepackage{upgreek}		% F�r nicht-kursive kleine griechischen Buchstaben

\usepackage{multirow}		% F�r mehrzeilige Felder in Tabellen

\usepackage{textcomp}		% F�r Sonderzeichen im normalen Text
							% (offensichtlich in tudreport schon eingebunden)
\usepackage[ngerman]{varioref}		% F�r vref
\usepackage{color}			% F�r farbigen Text
\usepackage{placeins}		% F�r \FloatBarrier
\usepackage{xspace}
\usepackage{icomma}			% Damit nach Dezimalkommas kein Abstand eingef�gt wird
							% (in math-Umgebungen)

\usepackage{cancel}			% Zum Wegstreichen von Gleichungstermen

\usepackage{array}			% F�r Zellentyp "m{}" in tabular-Umgebungen (Vertikal zentriert)

\usepackage{listings}		% Um formatierten Quellcode einzubinden
\usepackage{moreverb}		% F�r Umgebung "`verbatimtab"' (Verbatim mit Tabs)
\renewcommand{\verbatimtabsize}{4\relax}	% Standardtabweite in "`verbatimtab"'
											% ist 4 Zeichen


% Das Packet hyperref immer als letztes einbinden (nur bookmark darf danach kommen)!
%\usepackage[ps2pdf, colorlinks=false, pdfborder={0 0 0}]{hyperref}
%\usepackage[ps2pdf]{hyperref}	% F�r Verlinkungen im erzeugten pdf
\usepackage{hyperref}	% F�r Verlinkungen im erzeugten pdf
\usepackage{bookmark}
				% verwendete Pakete einbinden
% =================================================================================
% Definitionen f�r rtm / rtp
% =================================================================================
\ifrtm
	\newcommand{\SADAProf}{Prof. Dr.-Ing. Ulrich Konigorski}
  \newcommand{\SADAinstitut}{Institut f�r Automatisierungstechnik und Mechatronik\\
                             Fachgebiet Regelungstechnik und Mechatronik\\
                             \SADAProf}
  \newcommand{\SADAwebsite}{www.rtm.tu-darmstadt.de}
  \newcommand{\SADAtel}{06151/16-4167}
  \newcommand{\SADAlogo}{\includegraphics[width=3cm]{./common/IAT_rtm}}
\else
  \newcommand{\SADAProf}{Prof. Dr.-Ing. Dr. h.c. Rolf Isermann}
  \newcommand{\SADAinstitut}{Institut f�r Automatisierungstechnik\\
                             und Mechatronik\\
                             FG Regelungstechnik und Prozessautomatisierung\\
                             \SADAProf}
  \newcommand{\SADAwebsite}{http://www.rtm.tu-darmstadt.de/rtp.html}
  \newcommand{\SADAtel}{06151/16-5465}
  \newcommand{\SADAlogo}{\includegraphics[width=3cm]{./common/rtplogo}}
\fi
%\institution{\SADAinstitut}
% =================================================================================		
		
		
% =================================================================================
% Definitionen f�r die Erkl�rung, dass die Arbeit ohne Hilfe Dritter etc. erstellt wurde
% =================================================================================
% Die folgende Zeile deklariert die m�glichen Varianten der Erkl�rung, dass die
% Arbeit ohne Hilfe Dritter etc. erstellt wurde.
\def\ETIT{ETIT}\def\MBSA{MBSA}\def\MBDA{MBDA}
% =================================================================================


% =================================================================================
% Definitionen aus tudreport-Vorlage
% =================================================================================
\newif\ifTUDmargin
%\TUDmargintrue		% Sehr breiter rechter Rand
\TUDmarginfalse		% Schmaler (normaler) rechter Rand

\ifTUDmargin	% ggf. breiten Rand setzen
  \geometry{marginparsep=4.2mm,marginparwidth=28.64mm}
\fi

\newlength{\longtablewidth}
\setlength{\longtablewidth}{0.7\linewidth}
\addtolength{\longtablewidth}{-\marginparsep}
\addtolength{\longtablewidth}{-\marginparwidth}
% =================================================================================


% =================================================================================
% Farbige Deckseite, graue Balken auf allen anderen Seiten
% =================================================================================
\newif\ifOnlyColorFront
\OnlyColorFronttrue	    % Balken nur auf Deckblatt farbig
% \OnlyColorFrontfalse		% alle Balken farbig
% =================================================================================


% =================================================================================
% Befehle, die in scrreprt nicht exisiteren werden hier definiert, wenn diese
% Klasse verwendet wird.
% Auch die Seitenr�nder werden angepasst, so dass es grob wie mit tudreport
% aussieht.
% =================================================================================
\ifTUDdesign
	\title{\SADATitel}
	\subtitle{\SADAAutor}
	\subsubtitle{\SADATyp\ -- \SADAAbgabe  }%\\Betreuer: \SADABetreuer{} \SADABetreuerII{} \SADABetreuerIII{}}
	\setinstitutionlogo[height]{common/rtm_mit_schrift}
	\geometry{left=30mm, right=20mm, top=15mm, bottom=20mm}
\else
	\title{\SADATitel}
	\subtitle{}
	\author{\SADAAutor}	% scrreprt erwartet Autor
	
	\newcommand{\subsubtitle}[1]{}
	\newcommand{\settitlepicture}[1]{}
	\newcommand{\printpicturesize}{}
	\newcommand{\institution}[1]{}
    \pagestyle{headings}

  % Diese Pakete nur extra einbinden, wenn NICHT tudreport als Basis.
	\usepackage{amssymb}
	\usepackage{geometry}
	
% \settitlepicture{}	% Bild f�r Deckblatt
% \printpicturesize
\fi
% =================================================================================


% =================================================================================
% Anpassung Absatzformat
% =================================================================================
% Der Abstand der Zeilen betr�gt das 1,25-fache des Standard-Abstands von
% LaTeX. Da technische Arbeiten viele Formeln und Bilder enthalten, werden
% Abs�tze durch einen zus�tzlichen vertikalen Zwischenraum statt durch einen
% Einzug getrennt.
\linespread{1.25}
\setlength{\parindent}{0mm}
\setlength{\parskip}{1ex}
% =================================================================================


% =================================================================================
% Texte f�r die R�ckseite des Titelblatts vorgeben
% =================================================================================
\uppertitleback{}
\lowertitleback{}
\dedication{}
% =================================================================================


% =================================================================================
% Informationen (Meta-Daten) f�r pdf
% =================================================================================
\hypersetup{
	pdftitle = {\SADATitel},
	pdfsubject = {},
	pdfauthor = {\SADAAutor},
	pdfkeywords = {},
	pdfcreator = {},
	pdfproducer = {LaTeX with hyperref},
	pdfstartview = {Fit},
	pdfpagelayout = {SinglePage}
}
% =================================================================================
					% LaTeX-Einstellungen
% Dieses File dient zum definieren n�tzlicher Befehle.
% Es soll lediglich als Beispiel dienen, wie Befehle definiert werden, und welche Befehle n�tzlich sein k�nnen
%


% Inhalt
% ======
%	Allgemeine Abk�rzungen
%	Makros f�r Referenzen (Abbildungen, Zitate, ...)
%	Makros f�r Abbildungen
%	Makros f�r Einheiten, Exponenten
%	Makros f�r Formeln
%	Makros f�r Entwurf
%   Definitionen f�r Umgebungen



% Allgemeine Abk�rzungen
% ======================
	\newcommand{\bzw}{bzw.\@\xspace}
	\newcommand{\Bzw}{Bzw.\@\xspace}
	\newcommand{\bzgl}{bzgl.\@\xspace}
	\newcommand{\ca}{ca.\@\xspace}
	\newcommand{\evtl}{evtl.\@\xspace}
	\newcommand{\usw}{usw.\@\xspace}
	\newcommand{\etc}{etc.\@\xspace}
	\newcommand{\vgl}{vgl.\@\xspace}
	\newcommand{\Vgl}{Vgl.\@\xspace}
	\newcommand{\bspw}{bspw.\@\xspace}
	\newcommand{\Bspw}{Bspw.\@\xspace}
	\newcommand{\ggf}{ggf.\@\xspace}
	\newcommand{\Ggf}{Ggf.\@\xspace}



	\newcommand{\dah}{d.\thinspace{}h.\@\xspace}
	\newcommand{\Dah}{D.\thinspace{}h.\@\xspace}
	\newcommand{\iA}{i.\thinspace{}A.\@\xspace}
	\newcommand{\IA}{i.\thinspace{}A.\@\xspace}
	\newcommand{\ua}{u.\thinspace{}a.\@\xspace}
	\newcommand{\Ua}{U.\thinspace{}a.\@\xspace}
	\newcommand{\uU}{u.\thinspace{}U.\@\xspace}
	\newcommand{\UU}{u.\thinspace{}U.\@\xspace}
	\newcommand{\zB}{z.\thinspace{}B.\@\xspace}
	\newcommand{\ZB}{Zum Beispiel\xspace}
	\newcommand{\zT}{z.\thinspace{}T.\@\xspace}
	\newcommand{\ZT}{Z.\thinspace{}T.\@\xspace}



% Makros f�r Referenzen (Abbildungen, Zitate, ...)
% ================================================

	% Referenzierung auf Abbildungen, Tabellen, etc. (Hyperref-f�hig)
	\newcommand{\figref}[1]{\hyperref[#1]{\figurename\ \ref*{#1}}}
	\newcommand{\tabref}[1]{\hyperref[#1]{\tablename\ \ref*{#1}}}
	\newcommand{\equref}[1]{\hyperref[#1]{Gl.~(\ref*{#1})}}
	\newcommand{\defref}[1]{\hyperref[#1]{Definition~\ref*{#1}}}
	\newcommand{\figvref}[1]{\hyperref[#1]{\figurename}\vref{#1}}
	\newcommand{\tabvref}[1]{\hyperref[#1]{\tablename}\vref{#1}}
	\newcommand{\equvref}[1]{\hyperref[#1]{Gl.~(\ref*{#1}) auf Seite \pageref*{#1}}}
	\newcommand{\pagerefh}[1]{\hyperref[#1]{Seite~\pageref*{#1}}}
	\newcommand{\secref}[1]{\hyperref[#1]{Abschnitt~\ref*{#1}}}
	\newcommand{\charef}[1]{\hyperref[#1]{Kapitel~\ref*{#1}}}
	\newcommand{\lstref}[1]{\hyperref[#1]{Listing~\ref*{#1}}}


	% Zitate mit Seitenangabe in Fu�note
%	\newcommand{\citep}[2]{\cite{#1}\footnote{Seite #2}}
%	\newcommand{\citepp}[2]{\cite{#1}\footnote{Seiten #2}}
	\newcommand{\citep}[2]{\cite{#1} (S. #2)}
	\newcommand{\citepp}[2]{\cite{#1} (S. #2)}
	
	
% Makros f�r Abbildungen
% ======================
	% zum Skalieren nach Ersetzen durch psfrag
	\newcommand{\incgraphicsw}[2]{\resizebox{#1}{!}{\includegraphics{#2}}}


% Textbausteine
% =============
	% Produktnamen
	\newcommand{\Matlab}{\textsc{Matlab}}
	\newcommand{\Matlabreg}{\textsc{Matlab}\textsuperscript{\tiny \textregistered}}
	\newcommand{\MatSim}{\textsc{Matlab/Simulink}}
	\newcommand{\Simulink}{\textsc{Simulink}}
	\newcommand{\Simulinkreg}{\textsc{Simulink}\textsuperscript{\tiny \textregistered}}



% Makros f�r Einheiten, Exponenten
% ================================

	\newcommand{\unit}[1] { \ensuremath{\mathrm{#1}}}
	
	% Wert mit Einheit (mit kleinem Leerzeichen dazwischen), aus Text- UND Math-Modus
	\newcommand{\valunit}[2]{\ensuremath{#1\,\mrm{#2}}}


	% "�C", im Text- oder Mathe-Modus
	\newcommand{\degC}{
		\ifmmode
			^\circ \mrm{C}%
		\else
			\textdegree C%
		\fi}

	\newcommand{\degree}{
		\ifmmode
			^\circ%
		\else
			\textdegree%
		\fi}
	
	% F�r Exponentenschreibweise ( Anwendung: 123\E{3} )
	\newcommand{\E}[1]{ \ensuremath{\cdot 10^{#1}} }
	
	\newcommand{\eexp}[1]{ \mathrm{e}^{#1} }
	\newcommand{\iu}{ \mathrm{j} }

	\newcommand{\todots}{ ,\,\hdots,\, }


% Makros f�r Formeln
% ==================

    \newcommand{\mat}[1]{{\ensuremath{\boldsymbol{\mathrm{#1}}}}}

	\newcommand{\AP} { \mathrm{AP} }
	\newcommand{\doti} {(i)^\cdot}

	% Definition f�r Vektor und Matizen
	\newcommand{\ve}[1]{\ensuremath{\mathbf{#1}}}
	\newcommand{\ma}[1]{\ensuremath{\mathbf{#1}}}

	% Definition f�r Vektor und Matizen
	\newcommand{\ves}[1]{\ensuremath{\boldsymbol{#1}}}
	\newcommand{\mas}[1]{\ensuremath{\boldsymbol{#1}}}
	
	
	\newcommand{\inprod}[2]{\langle #1,\,#2 \rangle}
	
	\newcommand{\ul}[1]{\underline{#1}}

	% gerades "d" (z.B. f�r Integral)
	\newcommand{\ud} { \mathrm{d} }
	
	% normaler Text in Formeln
	\newcommand{\tn}[1] { \textnormal{#1} }
	
	% nicht-kursive Schrift in Formeln
	\newcommand{\mrm}[1] { \mathrm{#1}}
	
	% gerades "T" f�r Transponiert
	\newcommand{\transp}{\mathrm{T}}
	
	% gerades "rg"
	\newcommand{\rang}[1]{\mathrm{rg}(#1)}

	% F�r geklammerte Ausdr�cke mit Index (Subscript)
	% (einmal mit kursiven Index, einmal mit geradem Index)
	\newcommand{\grpsb}[2] { \left( #1 \right)_{#2} }
	\newcommand{\grprsb}[2] { \left( #1 \right)_{\mathrm{#2}} }

	% Ableitungen und Integrale
		% "normale" Ableitung (mit geraden "d"s)
		\newcommand{\normd}[2] { \frac{\mathrm{d} #1 }{\mathrm{d} #2 } }
		\newcommand{\normdat}[3] { \left. \frac{\mathrm{d} #1 }{\mathrm{d} #2 } \right|_{#3} }
	
		% Materielle Ableitung
		\newcommand{\matd}[2] { \frac{\mathrm{D} #1 }{\mathrm{D} #2 } }
		\newcommand{\matdat}[3] { \left. \frac{\mathrm{D} #1 }{\mathrm{D} #2 } \right|_{#3} }
	
		% Partielle Ableitung
		\newcommand{\partiald}[2] { \frac{\partial #1 }{\partial #2 } }
		\newcommand{\partialdat}[3] { \left. \frac{\partial #1 }{\partial #2 } \right|_{#3} }
	
	
	% Transformationen
	\newcommand{\FT}[1] { \mathfrak{F} \left\{ #1 \right\} }
	\newcommand{\FTabs}[1]{\left| \mathfrak{F} \left\{ #1 \right\} \right|}
	\newcommand{\IFT}[1] { \mathfrak{F}^{-1} \left\{ #1 \right\} }
	\newcommand{\DFT}[1]{\mathrm{DFT} \left\{ #1 \right\}}
	\newcommand{\DFTabs}[1]{\left| \mathrm{DFT} \left\{ #1 \right\} \right|}

	\newcommand{\Laplace}[1]{\mathfrak{L}\left (#1\right)} % L-Transformation
	\newcommand{\InvLaplace}[1]{\mathfrak{L^{-1}}\left (#1\right )} % L-Transformation, invers
	\newcommand{\invtrans}{\quad\bullet\!\!-\!\!\!-\!\!\circ\quad}
	\newcommand{\trans}{\quad\circ\!\!-\!\!\!-\!\!\bullet\quad}


	\newcommand{\mlfct}[1]{{\tt #1}}
	\newcommand{\mlvar}[1]{{\tt #1}}


	% Manche textcomp-Zeichen funktionieren mit dem TU-Design nicht, diese k�nnen dann mit diesem
	% Befehl gesetzt werden.
	\newcommand{\textcompstdfont}[1]{{\fontfamily{cmr} \fontseries{m} \fontshape{n} \selectfont #1}}
	


% =================================================================================
% Defintionen f�r Mathe-Umgebungen
% =================================================================================
	
\newtheorem{theorem}{Satz}
\newtheorem{lemma}[theorem]{Lemma}	% Selber Z�hler wie theorem
\newtheorem{definition}{Definition}
% =================================================================================


% =================================================================================
% Defintionen f�r Beispiel-Umgebung
% =================================================================================
\newcounter{examplenumber}[chapter]      % Neuer Counter bspnummer nummeriert nach Kapitel
\def\theexamplenumber{\thechapter.\arabic{examplenumber}}

\newenvironment{example}[1][]
{\vskip 3\parskip plus 1pt minus 1pt \refstepcounter{examplenumber}
\vspace{.3cm} \begin{addmargin}[1cm]{0cm} \noindent \textbf{Beispiel \theexamplenumber}: \textbf{#1}\par}
{\end{addmargin} \par \vspace{.3cm}}

% Alternative, einfachere Beispielumgebung:
% \newtheorem{example}{Beispiel}
% =================================================================================




% =================================================================================
% Definitionen f�r Listingsumgebung
% =================================================================================

\lstloadlanguages{Matlab}

\lstdefinestyle{Matlab_colored_smallfont}
{
	language = Matlab,
	tabsize = 4,
	framesep = 3mm,
	frame=tb,
	classoffset = 0,	
	basicstyle = \footnotesize\ttfamily,
	keywordstyle = \bfseries\color[rgb]{0,0,1},
	commentstyle = \itshape\color[rgb]{0.133,0.545,0.133},
	stringstyle = \color[rgb]{0.627,0.126,0.941},
	extendedchars = true,
	breaklines = true,
	prebreak = \textrightarrow,
	postbreak = \textleftarrow,
	%escapeinside = {(*@}{@*)},
	%moredelim = [s][\itshape\color[rgb]{0.5,0.5,0.5}]{[.}{.]},
	numbers = left,
	numberstyle = \tiny,
	stepnumber = 5
}

\lstdefinestyle{Matlab_colored}
{
	language = Matlab,
	tabsize = 4,
	framesep = 3mm,
	frame=tb,
	classoffset = 0,	
	basicstyle = \ttfamily,
	keywordstyle = \bfseries\color[rgb]{0,0,1},
	commentstyle = \itshape\color[rgb]{0.133,0.545,0.133},
	stringstyle = \color[rgb]{0.627,0.126,0.941},
	extendedchars = true,
	breaklines = true,
	prebreak = \textrightarrow,
	postbreak = \textleftarrow,
	%escapeinside = {(*@}{@*)},
	%moredelim = [s][\itshape\color[rgb]{0.5,0.5,0.5}]{[.}{.]},
	numbers = left,
	numberstyle = \tiny,
	stepnumber = 5
}


\lstdefinestyle{C_colored_smallfont}
{
	language=C,
	tabsize = 4,
	framesep = 3mm,
	frame=tb,	
	classoffset = 0,	
	basicstyle = \footnotesize\ttfamily,
	keywordstyle = \bfseries\color[rgb]{0,0,1},
	commentstyle = \itshape\color[rgb]{0.133,0.545,0.133},
	stringstyle = \color[rgb]{0.627,0.126,0.941},
	extendedchars = true,
	breaklines = true,
	prebreak = \textrightarrow,
	postbreak = \textleftarrow,
	%escapeinside = {(*@}{@*)},
	%moredelim = [s][\itshape\color[rgb]{0.5,0.5,0.5}]{[.}{.]},
	numbers = left,
	numberstyle = \tiny,
	stepnumber = 5
}

\lstdefinestyle{C_colored}
{
	language=C,
	tabsize = 4,
	framesep = 3mm,
	frame=tb,
	classoffset = 0,	
	basicstyle = \ttfamily,
	keywordstyle = \bfseries\color[rgb]{0,0,1},
	commentstyle = \itshape\color[rgb]{0.133,0.545,0.133},
	stringstyle = \color[rgb]{0.627,0.126,0.941},
	extendedchars = true,
	breaklines = true,
	prebreak = \textrightarrow,
	postbreak = \textleftarrow,
	%escapeinside = {(*@}{@*)},
	%moredelim = [s][\itshape\color[rgb]{0.5,0.5,0.5}]{[.}{.]},
	numbers = left,
	numberstyle = \tiny,
	stepnumber = 5
}
		% oft verwendete Befehle
% =================================================================================


% =================================================================================
% Hier beginnt das eigentliche Dokument
% =================================================================================

\begin{document}
%\selectlanguage{ngerman}
\maketitle

% Die Farbe der Identit�tsleiste wird auf Grau umgestellt, damit nicht alle Seiten
% farbig gedruckt werden m�ssen
\ifTUDdesign
	\ifOnlyColorFront	% ggf. nachfolgende Balken andere Farbe zuweisen
		\makeatletter 	% ben�tigt, um die @-Befehle auszuf�hren
    \renewcommand{\@TUD@largerulecolor}{\color{tud0b}}% am besten gleiche Farbe wie in der ersten Zeile und die Zahl durch die 0 ersetzen, dann hat das Grau die richtige Intensit�t
    \makeatother
	\fi
\fi


\pagenumbering{roman}	% Bis zum Hauptteil werden r�mische Seitenzahlen verwendet

% =================================================================================
% Spezielle Seiten f�r studentische Arbeiten
% =================================================================================
\cleardoublepage
%\section*{Aufgabenstellung}
F�r schriftliche Arbeiten (Pro-/Projektseminar, Studien-, Bachelor-, Master-, Diplomarbeiten, \etc) soll Studierenden ein \LaTeX-Dokument zur Verf�gung gestellt werden, das die Vorgaben aus den \emph{Richtlinien zur Anfertigung von Studien- und Diplomarbeiten}~\cite{Richtlinien} umsetzt. Die Dokumentation soll die Funktionen des Dokumentes beschreiben und Hinweise zu ihrer Anwendung geben.

Grundlage ist die \verb|tudreport|-Klasse. Die damit erstellten Arbeiten m�ssen sowohl zum Ausdrucken geeignet sein als auch f�r die Bildschirmdarstellung und die elektronische Archivierung als PDF-Datei.

\vspace{0.5cm}
\begin{tabular}{ll}
Beginn: & \SADABegin \\
Ende:   & \SADAAbgabe \\
Seminar:& \SADASeminar
\end{tabular}

\vspace{1cm}

\begin{tabular}{ll}
\rule{7cm}{0.4pt} \hspace{1cm} & \rule{7cm}{0.4pt} \\
\SADAProf & \SADABetreuer\\
 &\SADABetreuerII\\
 &\SADABetreuerIII
\end{tabular}

\vfill
{\renewcommand{\baselinestretch}{1} % f�r diesen Abschnitt einfacher Zeilenabstand
\normalsize % anwenden des Zeilenabstandes
\begin{minipage}{0.8\textwidth}
	Technische Universit�t Darmstadt\\
	\SADAinstitut\\[0.5cm]
%
	Landgraf-Georg-Stra�e 4\\
	64283 Darmstadt\\
	Telefon \SADAtel\\
	\SADAwebsite
\end{minipage}
\begin{minipage}{0.2\textwidth}
\flushright  % rechtsb�ndig
\ \\[2.7cm]
\SADAlogo\;
\end{minipage}}



%\cleardoublepage
%\ \\[3cm]	% Diese Zeile erzeugt einen Abstand von 4cm zur ersten Zeile, die nur ein Leerzeichen
			% enth�lt

\ifx\SADAVarianteErklaerung\ETIT
	\section*{Erkl�rung}
	\noindent
	Hiermit versichere ich, dass ich die vorliegende Arbeit ohne Hilfe Dritter und nur mit den angegebenen Quellen und Hilfsmitteln angefertigt habe. Alle Stellen, die aus den Quellen entnommen wurden, sind als solche kenntlich gemacht. Diese Arbeit hat in gleicher oder �hnlicher Form noch keiner Pr�fungsbeh�rde vorgelegen.\vspace*{20mm} \\
	\noindent
	\begin{tabular}{ll}
		\SADAStadt, den \SADAAbgabe	\hspace{1cm}	& \rule{0.4\textwidth}{0.4pt}\\
										            & \SADAAutor
	\end{tabular}
	


\else\ifx\SADAVarianteErklaerung\MBDA
	\section*{Erkl�rungen}
	\noindent
	Hiermit erkl�re ich an Eides statt, dass ich die vorliegende \SADATyp\ mit dem Titel\ \glqq\SADATitel\grqq\ selb\-st�ndig und ohne fremde Hilfe verfasst, andere als die angegebenen Quellen und Hilfsmittel nicht benutzt und die aus anderen	Quellen entnommenen Stellen als solche gekennzeichnet habe.\\
	Diese Arbeit hat in gleicher oder �hnlicher Form noch keiner Pr�fungsbeh�rde vorgelegen.\vspace*{20mm} \\
	\noindent
	\begin{tabular}{ll}
		\SADAStadt, den \SADAAbgabe	\hspace{1cm}	& \rule{0.4\textwidth}{0.4pt}\\
										& \SADAAutor
	\end{tabular}
	
	
	\vspace{40mm}
	\noindent
	Ich bin damit einverstanden, dass die TU Darmstadt das Urheberrecht an meiner \SADATyp\ zu wissenschaftlichen Zwecken nutzen kann.\vspace*{20mm} \\
	\noindent
	\begin{tabular}{ll}
		\SADAStadt, den \SADAAbgabe	\hspace{1cm}	& \rule{0.4\textwidth}{0.4pt}\\
										& \SADAAutor
	\end{tabular}

	{\huge Hier fehlt noch was!}
	
	
\else\ifx\SADAVarianteErklaerung\MBSA
	\section*{Erkl�rungen}
	\noindent
	Hiermit erkl�re ich an Eides statt, dass ich die vorliegende \SADATyp\ mit dem Titel\ \glqq\SADATitel\grqq\ selb\-st�ndig und ohne fremde Hilfe verfasst, andere als die angegebenen Quellen und Hilfsmittel nicht benutzt und die aus anderen	Quellen entnommenen Stellen als solche gekennzeichnet habe.\\
	Diese Arbeit hat in gleicher oder �hnlicher Form noch keiner Pr�fungsbeh�rde vorgelegen.\vspace*{20mm} \\
	\noindent
	\begin{tabular}{ll}
		\SADAStadt, den \SADAAbgabe	\hspace{1cm}	& \rule{0.4\textwidth}{0.4pt}\\
										& \SADAAutor
	\end{tabular}
	
	
	\vspace{40mm}
	\noindent
	Ich bin damit einverstanden, dass die TU Darmstadt das Urheberrecht an meiner \SADATyp\ zu wissenschaftlichen Zwecken nutzen kann.\vspace*{20mm} \\
	\noindent
	\begin{tabular}{ll}
		\SADAStadt, den \SADAAbgabe	\hspace{1cm}	& \rule{0.4\textwidth}{0.4pt}\\
										& \SADAAutor
	\end{tabular}


\else
	{\huge Unbekannte Variante der Erkl�rung!}

\fi\fi\fi





%
%\clearpage
%\section*{Kurzfassung}
Das \LaTeX-Dokument \verb|sada_tudreport| ist eine Vorlage f�r schriftliche Arbeiten (Proseminar-, Projektseminar-, Studien-, Bachelor-, Master- und Diplomarbeiten, \etc) am Institut f�r Automatisierungstechnik der TU Darmstadt. Das
Layout ist an die \emph{Richtlinien zur Anfertigung von Studien- und
Diplomarbeiten}~\cite{Richtlinien} angepasst und durch Modifikation der Klasse \verb|tudreport|
realisiert, so dass in der Arbeit die gewohnten \LaTeX-Befehle benutzt werden
k�nnen. Die vorliegende Anleitung beschreibt die Klasse und gibt grundlegende
Hinweise zum Verfassen wissenschaftlicher Arbeiten. Sie ist au�erdem ein
Beispiel f�r den Aufbau einer Studien-, Bachelor-, Master- bzw. Diplomarbeit.

\textbf{Schl�sselw�rter:} Studienarbeit, Bachelorarbeit, Masterarbeit, Diplomarbeit, Vorlage, \LaTeX-Klasse



\selectlanguage{english}
\section*{Abstract}
The \LaTeX\ document \verb|sada_tudreport| provides a template for student's research
reports and diploma theses (`` Proseminar-, Projektseminar-, Studien-, Bachelor-, Master- und Diplomarbeiten'') at the Institute of
Automatic Control, Technische Universit�t Darmstadt. The layout is adapted to
the \emph{``Richtlinien zur Anfertigung von Studien- und
Diplomarbeiten''}~\cite{Richtlinien} and is implemented by modification of the standard \verb|tudreport|
class, so that common \LaTeX\ commands can be used in the text. This manual
describes the class and dwells on general considerations on how to write
scientific reports. Additionally, it is an example for the structure of a
thesis.

\textbf{Keywords:} Research reports, diploma theses, template, \LaTeX\ class
\selectlanguage{ngerman} 
% =================================================================================

% =================================================================================
% Inhaltsverzeichnis
% =================================================================================
\cleardoublepage	% Auf einer leeren rechten Seite beginnen
\phantomsection		% Diese und die n�chste Zeile dient dazu, f�r das Inhalts-
					% verzeichnis einen Eintrag in das pdf-Inhaltsverzeichnis,
					% aber nicht in das normale Verzeichnis zu erzeugen.
%\pdfbookmark[0]{\contentsname}{pdf:toc}	
%\tableofcontents	% Inhaltsverzeichnis einf�gen
%\clearpage	% Sonst kommt nichts mehr auf die Seite
% =================================================================================


% =================================================================================
% Symbole und Abk�rzungen
% =================================================================================
% Nach dem Inhaltsverzeichnis kommt ein Verzeichnis der h�ufig verwendeten
% Symbole und Abk�rzungen. Dazu kann man das Paket 'nomencl' verwenden, oder man
% erstellt es von Hand.
%\chapter*{Symbole und Abk�rzungen}
\addcontentsline{toc}{chapter}{Symbole und Abk�rzungen} % erzeugt einen Eintrag im Inhaltsverzeichnis
%
\paragraph*{Lateinische Symbole und Formelzeichen}
\begin{tabularx}{\textwidth}{@{}l@{\qquad}X@{\quad}p{18mm}}
	Symbol & Beschreibung & Einheit\\ \midrule
  $I$ & Strom 			& \unit{A}\\
  $R$ & Widerstand 	& \unit{\Omega}\\
  $U$ & Spannung 		& \unit{V}\\
\end{tabularx}
%
\paragraph*{Griechische Symbole und Formelzeichen}
\begin{tabularx}{\textwidth}{@{}l@{\qquad}X@{\quad}p{18mm}}
	Symbol & Beschreibung & Einheit\\ \midrule
  $\mat{\Psi}$ & Datenmatrix\\
  $\sigma$     & Standardabweichung\\
  $\omega$     & Kreisfrequenz 	& \unit{s^{-1}}
\end{tabularx}
%
\paragraph*{Abk�rzungen}
\begin{tabularx}{\textwidth}{@{}l@{\qquad}X}
K�rzel & vollst�ndige Bezeichnung \\ \midrule
  Dgl. & Differentialgleichung\\
  LS   & Kleinste Quadrate (\emph{Least Squares})\\
  PRBS & Pseudo-Rausch-Bin�r-Signal (\emph{Pseudo Random Binary Signal})\\
  ZVF  & Zustandsvariablenfilter
\end{tabularx}
%
%\cleardoublepage

% =================================================================================
% Hauptteil
% =================================================================================
\pagenumbering{arabic}	% Hauptteil bekommt arabische Seitenzahlen % Titelseite, Aufgabenstellung, Erkl�rung, Abstract, Inhaltsverzeichnis, etc.
\maketitle

\setcounter{chapter}{2}

%Warum geht title/section nicht?
\section{Lienarisierung}
Im folgenden Abschnitt wird die Funktion zur Linearisierung des Doppelpendel-Systems um einen Arbeitspunkt, sowie ihre R�ckgabewerte an bestimmten Arbeitspunkten dokumentiert. Die Implementierung der Funktion ist in \autoref{Code_Linearisierung} aufgef�hrt.

Bevor das System linearisiert wird sind zwei Fragen zu kl�ren:
\begin{enumerate}
	\item Welche Arbeitspunkte sind sinnvoll?
	\item Was bedeutet es physikalisch, wenn $M_{\textrm{AP}}$ ungleich null ist?
\end{enumerate}
Die Antworten lauten wie folgt:
\begin{enumerate}
	\item Es ist nur sinnvoll das System in Arbeitspunkten zu linearisieren, in denen es sowohl vollst�ndig beobachtbar, als auch steuerbar ist.
	\item Bei der Gr��e $M_{\textrm{AP}}$ handelt es sich um den statischen Wert der Stellgr��e $M$ im Arbeitspunkt. Ist diese ungleich null muss der Motor das Moment $M_{\textrm{}}$ 
\end{enumerate}

\lstinputlisting[style=Matlab_colored,caption={Code der Linearisierungsfunktion}, label={Code_Linearisierung}]{Codes/linearisierung.m}

Die Linearisierung um die Arbeitspunkte
\begin{align*}
	\textbf{x}_{\textrm{AP}_1} &= \begin{bmatrix} 0 & 0 & 0 & 0\end{bmatrix}\\
	\textbf{x}_{\textrm{AP}_2} &= \begin{bmatrix} \pi & 0 & \pi & 0\end{bmatrix}\\
	\textbf{x}_{\textrm{AP}_3} &= \begin{bmatrix} \pi/2 & 0 & \pi & 0 \end{bmatrix}
\end{align*}
ergibt f�r die allgmeine Zustandsraumdarstellung:
\begin{align*}
	\textbf{\.{x}} &= \textbf{Ax + Bu}\\
	\textbf{y} &= \textbf{Cx + Du}
\end{align*}
die folgenden Systemmatrizen:
\begin{align*}
\textbf{A}_{\textrm{AP}_1} &= \begin{bmatrix} 0& 0 & 0& 0 \end{bmatrix}&
\textbf{B}_{\textrm{AP}_1} &= \begin{bmatrix} 0& 0 & 0 & 0 \end{bmatrix}\\
\end{align*}
\section{Vergleich der Linearisierten Modelle}

\section{Normalformen des Zustandsraummodelles}

\section{Untersuchung von Steuerbarkeit und Beobachtbarkeit}
\end{document}
